\documentclass{article}
\usepackage[T1]{fontenc}
\usepackage{lmodern}
\usepackage{graphicx}
\usepackage{amsmath,amssymb}
\usepackage[a4paper,margin=2cm]{geometry}
\usepackage[absolute,overlay]{textpos}
\setlength{\TPHorizModule}{1mm}
\setlength{\TPVertModule}{1mm}
\usepackage[indonesian]{babel}
\usepackage{hyperref}
\setlength{\emergencystretch}{2em}
% unit: Halaman 47

\begin{document}

% === Halaman 47 (python/native) ===

Encoding Pesan menjadi Titik di dalam Kurva

• Pesan yang akan dienkripsi dengan ECC harus dikonversi (encoding) menjadi titik 

di dalam kurva eliptik.

• Metode yang sederhana adalah memetakan setiap karakter ASCII dengan setiap 

titik pada kurva eliptik.

• Untuk 256 karakter ASCII, maka dibutuhkan kurva eliptik yang berisi minimal 256 

titik.

• Misalkan pesan M  = ‘ENCRYPT’, yang dalam nilai ASCII adalah ‘69’ ‘78’, ‘67’, ‘82’, 

‘89’, ‘80’, ‘84’. Setiap nilai ini dipetakan ke sebuah titik pada kurva eliptik.

• Namun metode ini kurang aman.

47

\end{document}
